\documentclass[conference]{IEEEtran}

\usepackage{amsmath}
\usepackage{booktabs}
\usepackage{tabularx}
\usepackage{url}
\usepackage{cite}
\usepackage[hidelinks]{hyperref}
\usepackage{xcolor}

\newcommand{\code}[1]{\texttt{#1}}
\newcommand{\pending}[1]{\textcolor{red}{\textbf{[#1]}}}
\newif\ifresultsready
\resultsreadyfalse

% Populate only from the one-shot human-final result artifact, then set \resultsreadytrue.
\newcommand{\FinalBinaryN}{\pending{N-BINARY}}
\newcommand{\FinalUnsafeN}{\pending{N-UNSAFE}}
\newcommand{\FinalNegativeN}{\pending{N-BOUNDED-NEGATIVE}}
\newcommand{\FinalIndeterminateN}{\pending{N-INDETERMINATE}}
\newcommand{\FinalAgreement}{\pending{RAW-AGREEMENT}}
\newcommand{\FinalKappa}{\pending{KAPPA}}
\newcommand{\PrimaryDelta}{\pending{FULL-MINUS-UNTYPED}}
\newcommand{\PrimaryLower}{\pending{PRIMARY-CI-LOWER}}
\newcommand{\PrimaryUpper}{\pending{PRIMARY-CI-UPPER}}
\newcommand{\FullAUPRC}{\pending{TBD}}
\newcommand{\HistAUPRC}{\pending{TBD}}
\newcommand{\SequenceAUPRC}{\pending{TBD}}
\newcommand{\CapabilityAUPRC}{\pending{TBD}}
\newcommand{\UntypedAUPRC}{\pending{TBD}}
\newcommand{\ActorRemovedAUPRC}{\pending{TBD}}
\newcommand{\ProjectBalancedAUPRC}{\pending{TBD}}
\newcommand{\FullAUROC}{\pending{TBD}}
\newcommand{\HistAUROC}{\pending{TBD}}
\newcommand{\SequenceAUROC}{\pending{TBD}}
\newcommand{\CapabilityAUROC}{\pending{TBD}}
\newcommand{\UntypedAUROC}{\pending{TBD}}
\newcommand{\ActorRemovedAUROC}{\pending{TBD}}
\newcommand{\ProjectBalancedAUROC}{\pending{TBD}}
\newcommand{\FinalDeferRate}{\pending{TBD}}
\newcommand{\FullWarnRecall}{\pending{TBD}}
\newcommand{\FullWarnRecallLower}{\pending{TBD}}
\newcommand{\FullWarnRecallUpper}{\pending{TBD}}
\newcommand{\FullWarnFPR}{\pending{TBD}}
\newcommand{\FullWarnFPRLower}{\pending{TBD}}
\newcommand{\FullWarnFPRUpper}{\pending{TBD}}
\newcommand{\FullWarnPrecision}{\pending{TBD}}
\newcommand{\FullWarnPrecisionLower}{\pending{TBD}}
\newcommand{\FullWarnPrecisionUpper}{\pending{TBD}}
\newcommand{\FullDeferLower}{\pending{TBD}}
\newcommand{\FullDeferUpper}{\pending{TBD}}
\newcommand{\FullNoWarnUnsafeRate}{\pending{TBD}}
\newcommand{\FullNoWarnUnsafeLower}{\pending{TBD}}
\newcommand{\FullNoWarnUnsafeUpper}{\pending{TBD}}
\newcommand{\PaperBranch}{\pending{METHOD-OR-MEASUREMENT}}
\InputIfFileExists{generated/prelabel_results_macros.tex}{}{\PackageError{AuthGuard}{Missing generated pre-label macros}{Run generate_prelabel_submission_assets.py}}
\InputIfFileExists{generated/final_results_macros.tex}{}{}

\title{AuthGuard-7702: Coverage-Audited Delegation-Context Analysis for
Pre-Authorization Screening}

\author{\IEEEauthorblockN{Anonymous Authors}}

\begin{document}
\maketitle

\ifresultsready\else
\begin{center}
\fbox{\parbox{0.94\columnwidth}{\centering\textbf{PRE-LABEL DRAFT---NOT FOR SUBMISSION.}
Final human-result macros are intentionally unresolved.}}
\end{center}
\fi

\begin{abstract}
EIP-7702 lets an externally owned account execute delegate-contract code in the account's
storage, balance, and identity context. This creates a security decision before authorization,
when behavioral history, reputation, and verified source may be unavailable. We present
AuthGuard-7702, a bytecode-only warning and triage study built around a Delegation-Context Risk
Graph (DCRG). DCRG connects reachable sensitive capabilities to recognized guard evidence and
retains unresolved control flow as an explicit coverage state. A bounded extractor combines
conservative state widening with jump-fenced Solidity-metadata recognition. On
\CanonicalRuntimes{} unique
runtimes, these corrections increase complete bounded-analysis coverage from 31.1\% to 63.8\%
without converting a previously complete runtime to partial coverage, while inherited-label
AUPRC remains statistically unchanged. We preregister a single representation endpoint on a
score-blind post-cutoff sample, retrain after conservative research-family holds, and separately
evaluate three new legitimate projects. \ifresultsready On \FinalBinaryN{} adjudicated binary
judgments, full DCRG reaches \FullAUPRC{} AUPRC and differs from its type-erased ablation by
\PrimaryDelta{} (95\% paired dependence-cluster interval [\PrimaryLower{}, \PrimaryUpper{}]).
\else Final untouched-label performance will be inserted only from the locked one-shot
evaluation.\fi The interface emits \code{WARN}, \code{NO\_MODEL\_WARNING}, or \code{DEFER}; it
does not infer intent or certify authorization safety.
\end{abstract}

\begin{IEEEkeywords}
EIP-7702, Ethereum, smart-account security, bytecode analysis, selective prediction
\end{IEEEkeywords}

\section{Introduction}

EIP-7702 changes the Ethereum account model by allowing an externally owned account (EOA) to
install a delegation indicator that resolves to contract code~\cite{eip7702}. The referenced code
then executes in the EOA's context. This enables batching, sponsorship, and scoped authorization,
but a deficient delegate can expose the authorizing account to replay, unexpected calls,
initialization takeover, storage collision, or broad asset control~\cite{eip7702}.

Existing EIP-7702 security studies establish that the threat is real. Huang et al. combine
large-scale historical transaction filtering, Gigahorse decompilation, and cross-contract rules to
identify EOA-targeted, contract-targeted, and composite attacks~\cite{huang2026darkside}. Qi et al.
study persistent delegation phishing, activation pathways, and ecosystem concentration
~\cite{qi2025phishing}. We ask a narrower operational question: what can be learned from a proposed
delegate's runtime bytecode \emph{before} authorization, without requiring behavioral history?

This decision boundary creates two methodological traps. First, the inherited malicious class was
constructed using static-analysis behavior, so a bytecode model can appear strong by reconstructing
the labeling rule. Second, exact-runtime separation does not prevent project or deployer dependence.
Our protocol therefore separates development labels from an untouched post-cutoff human study,
applies deliberately overinclusive research-family holds, freezes predictions before annotation,
and reports incomplete analysis as deferral.

This study targets three contributions:
\begin{itemize}
  \item a leakage-reduced, temporally separated evaluation resource for bytecode-only EIP-7702
  screening, retaining indeterminate and non-screenable outcomes;
  \item a coverage-audited, guard-aware Delegation-Context representation that relates reachable
  sensitive capabilities to authorization evidence; and
  \item a warning/no-model-warning/defer evaluation over post-cutoff signer/delegate pairs,
  conservative family holds, strong learned and classical baselines, and separately frozen
  legitimate-project controls.
\end{itemize}

These are bounded claims. DCRG is represented by aggregate features consumed by XGBoost; it is not
a learned graph neural network. \code{NO\_MODEL\_WARNING} is not ``safe.'' Anonymous linkage
clusters are retraining controls, not verified project attribution.

\section{Background and Related Work}

\subsection{EIP-7702 security}

An authorization tuple identifies a chain, delegate address, nonce, and EOA signature. A valid
tuple writes a delegation indicator to the EOA's code field; code-executing operations resolve the
referenced delegate~\cite{eip7702}. The specification explicitly warns that delegates should bind
replay, value, gas, target, and calldata, and that initialization and storage migration are
security-critical.

Huang et al. are the closest direct comparator~\cite{huang2026darkside}. Their transaction analyzer
narrows historical activity, after which Gigahorse and rule matching identify attack semantics.
AuthGuard-7702 does not replace that richer forensic setting; it studies score-blind screening at an
earlier decision with bytecode as the required input. Qi et al. provide a complementary measurement
and phishing analysis~\cite{qi2025phishing}. Accordingly, we make no claim to the first EIP-7702
detector, dataset, attack taxonomy, or connection to ERC-4337.

\subsection{Bytecode analysis and selective prediction}

Opcode models, CFG/data-flow analysis, program graphs, and cross-contract analysis are established
smart-contract-security techniques~\cite{gigahorse,eth2vec}. PhishingHook compares 16 histogram,
vision, language, and vulnerability-model techniques for generic phishing-contract classification
from EVM bytecode~\cite{phishinghook}. It therefore rules out first bytecode-phishing and first
opcode-model-comparison claims. Because its labels and population differ from EIP-7702 delegate
review, we do not compare published accuracy across datasets; our histogram+n-gram and learned
sequence baselines are retrained under the same holds and labels as DCRG. Classification with a
reject option is likewise established~\cite{selectivenet}. Our contribution is not a generic graph
or rejection algorithm; it is the EIP-7702-specific representation, coverage contract, and
provenance-aware evaluation at authorization time.

Full static analyzers such as Gigahorse and Securify recover richer program and compliance
semantics~\cite{gigahorse,securify}; AuthGuard does not claim to replace them. Its bounded extractor
exists to expose a compact, locally reproducible evidence contract and an explicit escalation path.
Security-ML evaluation is also vulnerable to spatial and temporal dependence: TESSERACT shows how
such biases inflate malware-classification estimates~\cite{tesseract}. We apply the same principle
to exact runtimes, signer/deployer clusters, research families, and time, while acknowledging that
the resulting controls cannot prove independence.

\section{Task and Threat Model}

The input is the deployed runtime bytecode of a proposed delegate. When available in the real-context
study, the authorizing EOA recovered from the authorization signature is neutral context. A positive
judgment (\code{UNSAFE}) requires a concrete authorization-relevant unsafe capability or condition.
The bounded negative label, \code{NO\_CONCRETE\_UNSAFE\_BEHAVIOR\_FOUND}, means only that the
reviewed evidence did not establish such behavior. \code{INDETERMINATE} and
\code{NOT\_BYTECODE\_SCREENABLE} are reported and excluded from binary metrics.

The adversary may deploy delegate bytecode with sensitive operations, misleading selectors,
unusual control flow, or appended bytes. We do not claim to resolve behavior dependent on unavailable
state, unresolved proxies, external contracts, off-chain intent, or transaction-specific effects.
Those cases should defer. Private-key compromise and UI deception independent of delegate semantics
are outside scope.

Formally, an item is $(B,A,C)$, where $B$ is runtime bytecode, $A$ is an optional recovered
authority, and $C$ is extraction coverage. The learned models return a score $s\in[0,1]$; held-out
canonical validation determines seed-specific warning threshold $t_k$. Neither $A$ nor post-cutoff
labels are used to fit a threshold. The binary estimand is ranking performance among adjudicated
\code{UNSAFE} and bounded-negative items. Indeterminate and non-screenable outcomes remain part of
the population description but not the binary estimand.

The defender acts before signing an authorization and may lack verified source, transaction
history, reputation, and downstream state. The adversary can choose bytecode, misleading selectors,
unusual dispatch, appended bytes, and open call/storage capabilities. We do not assume that opcode
presence proves exploitability or that a known project proves legitimacy. The intended response is
triage: a warning invites rejection or deeper inspection; deferral invites source audit, simulation,
allowlisting, or transaction-aware analysis. The model does not replace those controls.

We organize the study around three questions. \textbf{RQ1} asks whether the coverage repairs improve
analyzer validity without trading away visible gaps. \textbf{RQ2} asks whether typed guard evidence
improves post-cutoff ranking beyond an untyped guard representation and fair bytecode baselines.
\textbf{RQ3} asks how the frozen warning/deferral contract behaves on post-cutoff human judgments
and separately sourced legitimate projects.

\section{Delegation-Context Risk Graph}

\subsection{Evidence representation}

DCRG is an authority-relative evidence representation, not a claim of complete program semantics.
It records four conceptual entities: externally reachable entry regions, sensitive capabilities,
guard evidence, and coverage gaps. A relation records that a capability is reachable from an entry
region and which recognized guard evidence dominates or precedes the site. Table~\ref{tab:schema}
summarizes the fixed schema.

\begin{table}[t]
\centering
\caption{DCRG evidence schema. ``Unknown'' remains observable.}
\label{tab:schema}
\small
\begin{tabularx}{\columnwidth}{p{0.19\columnwidth}XX}
\toprule
Entity & Recorded evidence & Boundary when unresolved \\
\midrule
Entry & dispatcher targets, fallback, externally reachable regions & unreached target \\
Capability & calls, value movement, storage writes, approvals, creation, destruction & sensitive-site gap \\
Guard & self-call, signature, stored/fixed authority, EntryPoint, caller supplied, \code{tx.origin} & untyped/no recognized guard \\
Coverage & explored functions, unresolved jumps and state limits & \code{PARTIAL} \\
\bottomrule
\end{tabularx}
\end{table}

Capabilities include value-moving calls, externally parameterized targets or calldata,
approvals/transfers, storage modification, contract creation, and self-destruction. A capability is
not automatically unsafe: the purpose of the graph is to retain its authorization context. Guard
types include self-call, recovered-signature evidence, stored authority, fixed-address comparison,
ERC-4337 EntryPoint, caller-supplied comparison, and \code{tx.origin}. When the authorization tuple
supplies signer $A$, a fixed-address comparison is additionally recorded as matching or differing
from $A$; transaction sender is not substituted for signer.

The classifier consumes a fixed 27-dimensional aggregate vector from this graph: function and
coverage counts, guarded and unguarded sensitive-site counts, per-guard-type counts,
authority-match/mismatch counts, fallback reachability, and unreached-site counts. Thus DCRG is a
typed representation with graph-derived aggregates, not a learned message-passing architecture.

\subsection{Bounded extraction and coverage}

The extractor begins with linear disassembly, dispatcher recovery, bounded CFG construction, and a
symbolic stack capped at 64 values. It explores at most 60,000 states and 96 states per program
counter. After eight visits, loop-varying non-control constants are conservatively widened while
valid jump destinations remain concrete. Unresolved control transfers and exploration limits are
recorded as gaps rather than silently discarded.

Raw bytecode creates a second failure mode: bytes in a compiler metadata suffix can resemble
\code{JUMPDEST}, \code{CALL}, or \code{SELFDESTRUCT}. A suspected Solidity CBOR suffix is excluded
from the opcode backstop only if it decodes to a recognized shape, begins on an instruction
boundary, follows a terminal predecessor, and contains no disassembled jump destination. If any
fence fails, the bytes remain in analysis. These checks are correctness repairs, not a claim that
metadata removal itself is novel.

The extractor emits \code{COMPLETE} only when all bounded obligations are discharged; otherwise it
emits \code{PARTIAL}. No missing relation becomes evidence of absence. Figure~\ref{fig:workflow}
shows how this coverage state remains visible through the operating decision.

\begin{figure}[t]
\centering
\small
\setlength{\fboxsep}{3pt}
\fbox{runtime + optional signer}
$\rightarrow$ \fbox{bounded CFG/DCRG}
$\rightarrow$ \fbox{score + coverage}\\[3pt]
$\rightarrow$ \fbox{seed consensus}
$\rightarrow$ \fbox{WARN / NO WARNING / DEFER}
\caption{The screening path. Partial analysis and seed disagreement can only warn or defer; they
cannot produce \code{NO\_MODEL\_WARNING}.}
\label{fig:workflow}
\end{figure}

\subsection{Scoring model and falsification variants}

The primary classifier is XGBoost over the fixed aggregate vector. The capability-only variant
removes guard context; the untyped variant retains guard presence but erases guard type; the
actor-removed variant removes EntryPoint and authority-relative features. This decomposition makes
the central claim falsifiable: the typed-guard contribution is supported only if full DCRG exceeds
the untyped variant on the untouched endpoint. A relational GNN, local typed motifs, and fixed
sequence/DCRG fusion were tested on development evidence but did not meet their preregistered or
paired-improvement gates and are retained as negative results.

\section{Evaluation Protocol}

\subsection{Populations and labels}

The canonical development corpus contains \CanonicalRows{} rows representing
\CanonicalRuntimes{} unique runtimes. Its
inherited target remains useful for training and implementation diagnostics but overlaps the static
evidence represented by DCRG. A separate 150-item provisional proxy informed method selection and is
therefore development data.

The confirmatory population is a score-blind sample of 150 post-cutoff exact-runtime families. One
predeclared development overlap is excluded from scoring. A disjoint 150-item reserve is
cryptographically locked and cannot be inspected for model selection. Two independent primary
reviewers label every item, and one separate adjudicator resolves each disagreement. Model scores,
source labels, provisional labels, and other primary judgments are hidden.

\begin{table*}[t]
\centering
\caption{Study populations and their allowed role. ``Frozen'' means content and identifiers are
hashed before the corresponding labels or scores are opened.}
\label{tab:populations}
\small
\begin{tabularx}{0.96\textwidth}{lrrXX}
\toprule
Population & Items & Unique runtimes & Label/evidence source & Allowed role \\
\midrule
Canonical development & \CanonicalRows & \CanonicalRuntimes & inherited static source rule & training and engineering diagnostics \\
Gold-Test proxy & 150 & 150 families & provisional, method-visible review & method selection only \\
Post-cutoff primary & 150 & 150 families & independent dual review; pending & confirmatory, 149 score-eligible \\
Replication reserve & 150 & 150 families & unopened & replacement/replication under locked rule \\
New legitimate projects & 3 & 5 controls & provenance-backed operating cases & descriptive warning/defer behavior \\
\bottomrule
\end{tabularx}
\end{table*}

\subsection{Post-cutoff construction}

The Ethereum collector scans a fixed checkpoint through block 24,641,536: \SnapshotBlocks{} blocks,
\SnapshotTypeFourTxs{} type-4 transactions, and \SnapshotAuthorizations{} authorization entries,
with no block-scan RPC error.
Zero-address revocations are excluded. We recover the signer from each authorization tuple rather
than using transaction sender; the two differ for \AuthoritySenderDifferences{} of
\SnapshotUsablePairs{} usable signer/delegate pairs. Runtime
code is read at first observed block. This is end-of-block state, so same-block changes remain a
limitation.

Among \SnapshotNonzeroDelegates{} distinct nonzero delegate addresses, \SnapshotUsablePairs{}
pairs have historical runtime code and yield 673
exact-runtime families. Exact hash and opcode-4-gram similarity against every canonical runtime
leave \SnapshotUnseenFamilies{} eligible unseen families. A deterministic sampler reads no label or model column and
selects at most one address per exact-runtime family. The primary and reserve use different frozen
seeds and share neither item nor exact-runtime family.

\subsection{Human judgment protocol}

Two qualified primary reviewers independently inspect every primary item using the fixed taxonomy:
\code{UNSAFE}, \code{NO\_CONCRETE\_UNSAFE\_BEHAVIOR\_FOUND}, \code{INDETERMINATE}, and
\code{NOT\_BYTECODE\_SCREENABLE}. Evidence packets contain runtime, recovered signer context,
decompilation and provenance aids, but recursively exclude all model scores, source-rule labels,
provisional labels, and another primary judgment. Reviewers record the concrete capability,
authority/caller condition, evidence boundary, and confidence. A distinct third reviewer
adjudicates every primary disagreement. We report pre-adjudication marginals, raw agreement,
Cohen's $\kappa$ when defined for the fixed pair, disagreement count, adjudication count, and final
class counts. The evaluator refuses a partial release or a release not exactly bound to all 150
manifest IDs.

\subsection{Leakage reduction and models}

Repeated recovered authorities, EOA deployers, verified families, and overinclusive bytecode
similarity links define conservative research clusters. These links produce mandatory canonical
family holds; they do not prove brand identity or complete independence. Three fixed seeds train a
learned sequence+dense baseline, normalized opcode histogram plus hashed 4-gram XGBoost,
capability-only DCRG, untyped guards, actor-removed DCRG, full DCRG, and a separately named
project-balanced DCRG.

The 150 primary items form \DependenceClusters{} dependence clusters;
\DependenceMultiClusters{} are multi-item clusters, \DependenceMultiItems{} items participate
in a multi-item cluster, and the maximum size is four. Must-links use repeated recovered-authority
EOAs and deployer EOAs; shared contract factories alone are not treated as ownership. A separate
score-blind audit labels one item as a confirmed development overlap, assigns 148 anonymous items
to conservative non-attribution clusters, and predeclares the overlap as excluded. Similarity links
generate \HeldCanonicalFamilies{} mandatory canonical-family holds. This intentionally
overinclusive procedure supports
``leakage-reduced,'' not ``leakage-free.''

After holds, each seed fits on \FrozenTrainRows{} training and \FrozenValidationRows{} validation
rows. The sequence baseline uses
chunk attention over opcode tokens plus a dense structural branch. The classical baseline combines
a normalized 225-bin opcode histogram with 512 hashed opcode 4-grams and XGBoost. DCRG variants use
the same fixed seed set. The project-balanced variant additionally assigns total training weight
eight to each of eight provenance-backed legitimate development projects; it is named separately
and is not substituted for the primary representation comparison. Nominal 5\% FPR thresholds and
temperature calibration use held-out canonical validation after all holds. Retraining freezes 21
\FrozenCheckpoints{} checkpoints and \FrozenPredictionRows{} primary score rows before annotation.

The single confirmatory endpoint is mean-seed AUPRC(full DCRG) minus AUPRC(untyped guards). Its 95\%
interval uses 10,000 paired seed-aware dependence-cluster bootstrap replicates. Baseline comparisons,
AUROC, calibration, operating metrics, coverage, and error analyses are secondary or descriptive.
Prevalence is always reported beside AUPRC.

For every bootstrap replicate, one dependence-cluster multiset is drawn and reused for all models
and all three seeds. Candidate-minus-baseline deltas are computed per seed and averaged within the
replicate before percentile bounds are taken. This avoids both item-level pseudoreplication and the
invalid practice of averaging independently computed confidence-interval endpoints. The primary
typed-guard hypothesis is supported only when the full-minus-untyped 95\% interval has lower bound
above zero. Full-minus-sequence and full-minus-histogram comparisons receive their own intervals;
no global baseline-superiority statement is inferred from the primary ablation.

\subsection{Decision contract and external controls}

A seed consensus at the validation-derived nominal 5\% FPR threshold emits \code{WARN} when at
least two of three seeds warn. Exactly one warning vote emits \code{DEFER}. With zero warning votes,
\code{PARTIAL} coverage or missing authority emits \code{DEFER}; only complete covered evidence can
emit \code{NO\_MODEL\_WARNING}. A failed development experiment showed that low-score cases could
not be certified as low risk, so no low-risk tier is retained.
This consensus contract is a label-blind supplemental analysis frozen while the post-cutoff
annotation count was zero; it does not change the preregistered AUPRC endpoint.

The separate pre-scoring external-control protocol contains Tangem, Startale, and Rainbow. It
distinguishes a new project from a new runtime and a new implementation lineage. Because all three
are legitimate operating cases and the sample is small, we report only project-level decisions and
never AUPRC, AUROC, significance, or a population false-positive rate.

For final operating analysis, warning recall, observed warning FPR, precision, defer rate, and the
observed unsafe fraction inside \code{NO\_MODEL\_WARNING} receive Wilson intervals. The latter is a
failure audit, not an estimated safety guarantee. Calibration, coverage, and an error taxonomy by
guard type, proxy pattern, bytecode size, and project novelty remain descriptive and cannot trigger
model revision.

\section{Results}

\subsection{RQ1: Coverage repair}

Complete bounded-analysis coverage increases from \BaseCompleteRuntimes{}/\CanonicalRuntimes{}
(31.1\%) to \FinalCompleteRuntimes{}/\CanonicalRuntimes{} (63.8\%) unique
runtimes, with no complete-to-partial regression. Mean inherited-label DCRG AUPRC changes from
0.95409 to 0.95375 ($-0.00034$, paired 95\% CI $[-0.00617,0.00423]$). Thus the repair improves
analysis validity without being selected for performance.

\begin{table}[t]
\centering
\caption{Label-free bounded-analysis coverage.}
\label{tab:coverage}
\small
\begin{tabular}{lrr}
\toprule
Extractor & Complete rows & Complete runtimes \\
\midrule
Original bounded CFG & \BaseCompleteRows{} (30.6\%) & \BaseCompleteRuntimes{} (31.1\%) \\
Metadata + widening & \IntermediateCompleteRows{} (66.3\%) & \IntermediateCompleteRuntimes{} (65.1\%) \\
Jump-fenced final & \FinalCompleteRows{} (64.7\%) & \FinalCompleteRuntimes{} (63.8\%) \\
\bottomrule
\end{tabular}
\end{table}

The exact-metadata intermediate appears slightly higher because syntactically valid trailers can
still contain a reachable instruction boundary or jump target. The final fences intentionally
restore 21 runtimes to \code{PARTIAL}. Across \CanonicalRuntimes{} runtimes the final extractor
reports no analysis
error; widening is an over-approximation and residual unresolved transfers remain visible.

\subsection{Development diagnostics, not final evidence}

The provisional 150-item proxy contains 108 binary and 42 uncertain judgments and was inspected
during method development. On it, full DCRG reaches 0.93885 mean AUPRC; full-minus-untyped is
$+0.00427$ (95\% paired family CI $[0.00076,0.01023]$), full-minus-sequence is $+0.04615$
($[0.00449,0.09178]$), and full-minus-histogram is $+0.03154$
($[-0.01116,0.08209]$). These diagnostics justify completing the expensive independent study, but
cannot appear in the abstract or support a submission claim because the proxy influenced method
selection.

\subsection{Frozen post-cutoff execution}

Post-cutoff retraining produced \FrozenCheckpoints{} locked checkpoints and
\FrozenPredictionRows{} score-only prediction rows (\FrozenScoredItems{} items
$\times$ three seeds). All source, feature, hold-plan, checkpoint, and prediction hashes passed
before review was unlocked. Table~\ref{tab:decisions} reports the resulting label-free consensus
counts. Similar counts across variants are not evidence of equivalent accuracy: the items within
each outcome can differ.

\begin{table}[t]
\centering
\caption{Pre-label decisions on 149 eligible primary items.}
\label{tab:decisions}
\small
\begin{tabular}{lrrr}
\toprule
Model & Warn & No warning & Defer \\
\midrule
Sequence+dense & \SequencePreWarn & \SequencePreNoWarning & \SequencePreDefer \\
Histogram+4-gram & \HistogramPreWarn & \HistogramPreNoWarning & \HistogramPreDefer \\
Capability-only & \CapabilityPreWarn & \CapabilityPreNoWarning & \CapabilityPreDefer \\
Untyped guards & \UntypedPreWarn & \UntypedPreNoWarning & \UntypedPreDefer \\
Actor-removed & \ActorRemovedPreWarn & \ActorRemovedPreNoWarning & \ActorRemovedPreDefer \\
Full DCRG & \FullDCRGPreWarn & \FullDCRGPreNoWarning & \FullDCRGPreDefer \\
Project-balanced & \ProjectBalancedPreWarn & \ProjectBalancedPreNoWarning & \ProjectBalancedPreDefer \\
\bottomrule
\end{tabular}
\end{table}

\subsection{External legitimate controls}

The external principal model emits \ExternalWarn{} \code{WARN}, \ExternalNoWarning{}
\code{NO\_MODEL\_WARNING}, and \ExternalDefer{} \code{DEFER}. Tangem and Rainbow are complete
no-warning cases; Startale defers under partial
signer-relative coverage. Rainbow is a new project but shares a known framework lineage, while the
canonical-unseen subset has one no-warning and one defer. Table~\ref{tab:external} keeps these
distinctions explicit. This $n=3$ project result is descriptive and does not estimate a false-alert
rate.

\begin{table}[t]
\centering
\caption{Principal decisions on new legitimate projects.}
\label{tab:external}
\small
\begin{tabular}{lcccc}
\toprule
Project & New? & Unseen? & Cov. & Decision \\
\midrule
Tangem & yes & yes & complete & no warning \\
Startale & yes & yes & partial & defer \\
Rainbow & yes & no & complete & no warning \\
\bottomrule
\end{tabular}
\end{table}

\subsection{RQ2--RQ3: Untouched human endpoint}

\begin{table}[t]
\centering
\caption{Post-cutoff model results. Red entries remain locked until independent review.}
\label{tab:final-models}
\small
\begin{tabular}{lcc}
\toprule
Model & AUPRC & AUROC \\
\midrule
Sequence+dense & \SequenceAUPRC & \SequenceAUROC \\
Histogram+4-gram & \HistAUPRC & \HistAUROC \\
Capability-only & \CapabilityAUPRC & \CapabilityAUROC \\
Untyped guards & \UntypedAUPRC & \UntypedAUROC \\
Actor-removed & \ActorRemovedAUPRC & \ActorRemovedAUROC \\
Full DCRG & \FullAUPRC & \FullAUROC \\
Project-balanced & \ProjectBalancedAUPRC & \ProjectBalancedAUROC \\
\bottomrule
\end{tabular}
\end{table}

The frozen release contains \FinalUnsafeN{} \code{UNSAFE}, \FinalNegativeN{} bounded-negative,
and \FinalIndeterminateN{} indeterminate/non-screenable judgments; \FinalBinaryN{} enter binary
metrics. Raw primary agreement is \FinalAgreement{} and Cohen's $\kappa$ is \FinalKappa{}.
Full-minus-untyped AUPRC is \PrimaryDelta{} with 95\% paired interval
[\PrimaryLower{}, \PrimaryUpper{}]. For full DCRG, warning recall is \FullWarnRecall{}, observed
warning FPR is \FullWarnFPR{}, warning precision is \FullWarnPrecision{}, defer rate is
\FinalDeferRate{}, and the observed unsafe fraction within \code{NO\_MODEL\_WARNING} is
\FullNoWarnUnsafeRate{}.

\begin{table}[t]
\centering
\caption{Full-DCRG frozen operating outcomes. Intervals are Wilson 95\%.}
\label{tab:operating}
\small
\begin{tabular}{lcc}
\toprule
Metric & Estimate & 95\% interval \\
\midrule
Warning recall & \FullWarnRecall & [\FullWarnRecallLower, \FullWarnRecallUpper] \\
Observed warning FPR & \FullWarnFPR & [\FullWarnFPRLower, \FullWarnFPRUpper] \\
Warning precision & \FullWarnPrecision & [\FullWarnPrecisionLower, \FullWarnPrecisionUpper] \\
Defer rate & \FinalDeferRate & [\FullDeferLower, \FullDeferUpper] \\
Unsafe within no-warning & \FullNoWarnUnsafeRate & [\FullNoWarnUnsafeLower, \FullNoWarnUnsafeUpper] \\
\bottomrule
\end{tabular}
\end{table}

\ifresultsready
\subsection{Primary decision}
The artifact-selected paper branch is \PaperBranch{}. The corresponding claim wording and all
secondary comparisons remain governed by the evidence gates in the submission blueprint.
\else
No method-superiority conclusion is available before the independent release. The proxy-label
diagnostic is intentionally excluded from this results table because it informed method selection.
\fi

\section{Discussion}

\subsection{Interpreting the result branches}

The final result must be interpreted at three levels. First, the primary interval determines whether
typed guard categories add measurable ranking value beyond guard presence. Second, fair sequence and
classical comparisons determine whether any baseline-superiority wording is supported. Third,
coverage and operating decisions determine whether the representation is useful as bounded triage
rather than merely a ranking model.

If the primary interval is strictly positive, the method contribution is evidence that guard type
adds ranking value under the frozen post-cutoff design. It still does not imply that every guard
category helps: actor-removed and capability-only comparisons are separately reported. If the
interval crosses zero, the method-superiority hypothesis is rejected for this study. The paper then
centers the measurement contribution: how label construction, research-family dependence, coverage
correction, and explicit deferral change the evaluation of bytecode-only screening. The branch is
selected mechanically and a null result is not rewritten as equivalence.

\subsection{Role in a defense stack}

The screening point exists after a wallet learns the delegate address but before it signs the
authorization. A \code{WARN} can block or demand stronger confirmation. A \code{DEFER} indicates
that source audit, state-aware simulation, cross-contract analysis, reputation, or an allowlist is
needed. \code{NO\_MODEL\_WARNING} says only that the fixed three-seed criterion was not met under
complete bounded analysis. This division is particularly important because the EIP specification
notes that changing a delegation is security-critical and that future balance behavior cannot
always be inferred statically~\cite{eip7702}.

\subsection{Negative paths as evidence}

The negative paths matter. Sequence fusion did not reliably improve DCRG; the GNN met a preregistered
futility stop; typed local motifs did not add a paired gain; protocol-actor features were neutral on
development evidence; and the attempted low-risk tier failed. Flood-200 evaluates one synthetic
transformation only. First-\code{STOP} canonicalization was not semantics-preserving and is not a
defense.

\begin{table}[t]
\centering
\caption{Method paths tested before the final freeze.}
\label{tab:negative}
\small
\begin{tabularx}{\columnwidth}{lXX}
\toprule
Path & Evidence gate & Disposition \\
\midrule
Relational GNN & futility upper bound & stopped \\
Typed motifs & paired gain over aggregate DCRG & null \\
Sequence fusion & paired gain over DCRG & supporting only \\
Low-risk tier & bounded risk under shift & failed; removed \\
First-\code{STOP} normalization & semantic preservation & failed; removed \\
Flood-200 & one valid transformation & bounded robustness result \\
\bottomrule
\end{tabularx}
\end{table}

These failures narrow the claim. In particular, aggregate DCRG should not be described as a GNN;
deferral should not be described as a new selective-learning algorithm; and Flood-200 should not be
generalized to adversarial robustness.

\section{Limitations and Ethics}

The post-cutoff study is Ethereum-only and covers a bounded temporal snapshot. Human judgments are
evidence-bounded rather than ground truth. The analyzer is neither sound nor complete for the EVM.
Training labels still originate from a static source rule, although confirmatory labels and
predictions are separated. Conservative similarity holds may overexclude while still failing to
prove all project dependence absent. The external legitimate set is only three projects, one case is
partial, and one shares a known framework lineage. Latency measurements do not establish wallet
integration or deployment safety.

The historical runtime query observes end-of-block state rather than transaction-index pre-state.
Conservative family links can both over-hold unrelated code and miss semantic or organizational
relations not visible in public provenance. Shared-authority/deployer clusters address statistical
dependence but do not yield project attribution. The fixed 150-item sample may contain too few
bounded-negative or positive judgments for stable operating estimates; the protocol therefore
requires both classes and reports intervals rather than silently proceeding.

Training on inherited source-rule labels can teach the model to reconstruct that rule. Independent
post-cutoff labels prevent direct test-label circularity but cannot erase inductive bias from
training. Likewise, a legitimate control is a documented operating case, not proof that all its
behavior is safe. The three-project control set is intentionally reported project by project.

We analyze public bytecode, authorization records, project documentation, and audit links. A model
warning is not an accusation, and a no-warning result is not an endorsement of any named project.
Reviewer identities may remain pseudonymous in the artifact, but qualifications, independence,
calibration, conflicts, and the private identity mapping must be retained for research integrity.

\section{Reproducibility and Artifact}

The artifact contains the canonical benchmark, exact-runtime and family identifiers, frozen
post-cutoff and reserve manifests, public-provenance worklists, DCRG graphs/features, seven model
variants, 21 checkpoints, score-only predictions, dependence clusters, legitimate-control
evidence, and the annotation application. SHA-256 locks bind the canonical dataset, sampling code,
review manifest, dependence clusters, family-hold plan, training code, feature definitions,
checkpoints, and predictions. The final evaluator verifies these identities before reading labels.

The implementation uses Python 3.12.12, evmole 0.9.0, cbor2 5.7.1, XGBoost, and PyTorch. Fixed seeds
are 7702, 7703, and 7704. The full revision-v3 suite contains 318 passing tests, including source
locks, family isolation, prediction cardinality, dual-review/adjudication gates, label-free evidence
packets, bootstrap behavior, decision consensus, and manuscript-claim auditing. The supplied build
container produces a letter-paper IEEE PDF. Final LaTeX macros are generated only from the two
verified evaluation reports and record both input hashes, avoiding manual result transcription.

\section{Conclusion}

AuthGuard-7702 establishes an auditable experimental framework for bytecode-only EIP-7702
pre-authorization screening. \ifresultsready \pending{INSERT THE SUPPORTED METHOD OR MEASUREMENT
CONCLUSION ONLY.} \else Whether typed guard evidence provides an independent predictive advantage
remains intentionally unresolved until the preregistered human evaluation.\fi\ In all outcomes, the
system prioritizes warning or deeper review; it does not prove malicious intent, exploitability, or
safety.

\bibliographystyle{IEEEtran}
\bibliography{references}

\end{document}
